% GENERATED FILE. Edit canonical lessons, then run npm run generate:books.
% canonical-source-hash: fnv1a64:6a1369d97b7b2cea
% canonical-lessons: KA-C101-open, KA-C101-close, KA-C101-stop, KA-C101-call, KA-C101-urgent

\chapter{Open, Close, Stop, Call}
\label{ch:ka-open-close-stop-call}
\hlchaptermodality{101}

\begin{chapteropening}
\textbf{By the end of this chapter:} \emph{I can politely ask someone to open, close, stop or call, and say it is urgent.}

Complete the last lesson of chapter 101: I can politely ask someone to open, close, stop or call, and say it is urgent.
\end{chapteropening}

\section[\texorpdfstring{\kn{ತೆರೆಯಿರಿ} (tereyiri)}{ತೆರೆಯಿರಿ (tereyiri)}]{\kn{ತೆರೆಯಿರಿ} (tereyiri) — please open}

\label{lesson:KA-C101-open}



\begin{quote}
Before the new one: say the Kannada for a station, then the Kannada for near.
\end{quote}

\subsection*{You'll want to know: \kn{ತೆರೆಯಿರಿ}}
\textbf{\kn{ತೆರೆಯಿರಿ}} (\emph{tereyiri}) — \textquotedblleft{}please open\textquotedblright{}, the polite \textbf{‑\kn{ಇರಿ}} form, like \textbf{\kn{ಕಾಯಿರಿ}}, \textquotedblleft{}please wait\textquotedblright{}.

\textbf{\kn{ಬಾಗಿಲು} \kn{ತೆರೆಯಿರಿ}} (\emph{bāgilu tereyiri}) — \textquotedblleft{}please open the door\textquotedblright{}.

\begin{cousinweb}
\textbf{\kn{ತೆರೆಯಿರಿ}} (\emph{tereyiri}) is native Dravidian, of the root for \textquotedblleft{}opening\textquotedblright{}, and DEDR 3259 lists \emph{teṟe}, \textquotedblleft{}to open, uncover\textquotedblright{}, itself; the entry also has Tamil \emph{tiṟa}, Telugu \emph{teṟa} and Malayalam \emph{tuṟakka}.
\end{cousinweb}

\begin{grammarlens}[title={What you've built}]
A polite command on a known ending.
\end{grammarlens}

\subsection*{Your turn}
\begin{itemize}
\raggedright
  \item \emph{Say it:} \emph{tereyiri}
  \item \emph{Say it:} \emph{tereyiri}, once more
  \item \emph{Say it:} say \emph{bāgilu tereyiri}
  \item \emph{Recall:} say the Kannada for a station, then the Kannada for near, then say \emph{tereyiri} again
\end{itemize}

\begin{tcolorbox}[breakable,colback=teal!4,colframe=teal!35!black,title={Before you move on}]
What does \kn{ತೆರೆಯಿರಿ} mean? (\textquotedblleft{}Please open\textquotedblright{}.) Say it once more.
\end{tcolorbox}
\section[\texorpdfstring{\kn{ಮುಚ್ಚಿ} (mucci)}{ಮುಚ್ಚಿ (mucci)}]{\kn{ಮುಚ್ಚಿ} (mucci) — please close}

\label{lesson:KA-C101-close}



\begin{quote}
Before the new one: say the Kannada for near, then the Kannada for please open.
\end{quote}

\subsection*{You'll want to know: \kn{ಮುಚ್ಚಿ}}
\textbf{\kn{ಮುಚ್ಚಿ}} (\emph{mucci}) — \textquotedblleft{}please close\textquotedblright{}.

\textbf{\kn{ಬಾಗಿಲು} \kn{ಮುಚ್ಚಿ}} (\emph{bāgilu mucci}) — \textquotedblleft{}please close the door\textquotedblright{}.

\begin{cousinweb}
\textbf{\kn{ಮುಚ್ಚಿ}} (\emph{mucci}) is native Dravidian, from \emph{muccu}, \textquotedblleft{}to close, shut\textquotedblright{}, which DEDR 4915 lists itself; \emph{mucci} is its polite command; the entry also has Tamil \emph{muccu} and Telugu \emph{mūyu}.
\end{cousinweb}

\begin{grammarlens}[title={What you've built}]
The partner of open.
\end{grammarlens}

\subsection*{Your turn}
\begin{itemize}
\raggedright
  \item \emph{Say it:} \emph{mucci}
  \item \emph{Say it:} \emph{mucci}, once more
  \item \emph{Say it:} say \emph{bāgilu mucci}
  \item \emph{Recall:} say the Kannada for near, then the Kannada for please open, then say \emph{mucci} again
\end{itemize}

\begin{tcolorbox}[breakable,colback=teal!4,colframe=teal!35!black,title={Before you move on}]
What does \kn{ಮುಚ್ಚಿ} mean? (\textquotedblleft{}Please close\textquotedblright{}.) Say it once more.
\end{tcolorbox}
\section[\texorpdfstring{\kn{ನಿಲ್ಲಿಸಿ} (nillisi)}{ನಿಲ್ಲಿಸಿ (nillisi)}]{\kn{ನಿಲ್ಲಿಸಿ} (nillisi) — please stop}

\label{lesson:KA-C101-stop}



\begin{quote}
Before the new one: say the Kannada for please open, then the Kannada for please close.
\end{quote}

\subsection*{You'll want to know: \kn{ನಿಲ್ಲಿಸಿ}}
\textbf{\kn{ನಿಲ್ಲಿಸಿ}} (\emph{nillisi}) — \textquotedblleft{}please stop (it)\textquotedblright{}, said to a driver or about a machine. It comes from \textbf{\kn{ನಿಲ್ಲು}}, \textquotedblleft{}to stop, to stand\textquotedblright{}.

\begin{cousinweb}
\textbf{\kn{ನಿಲ್ಲಿಸಿ}} (\emph{nillisi}) is native Dravidian, from \emph{nil}, \textquotedblleft{}to stand, stop\textquotedblright{}, with \emph{-isu}, \textquotedblleft{}to make\textquotedblright{}: \emph{nillisu}, \textquotedblleft{}to make stop\textquotedblright{}; DEDR 3675 lists the root with Tamil \emph{nil}, Telugu \emph{nilucu} and Malayalam \emph{nilkka}.
\end{cousinweb}

\begin{grammarlens}[title={What you've built}]
The request on a bus or in a taxi.
\end{grammarlens}

\subsection*{Your turn}
\begin{itemize}
\raggedright
  \item \emph{Say it:} \emph{nillisi}
  \item \emph{Say it:} \emph{nillisi}, once more
  \item \emph{Say it:} say \emph{illi nillisi}
  \item \emph{Recall:} say the Kannada for please open, then the Kannada for please close, then say \emph{nillisi} again
\end{itemize}

\begin{tcolorbox}[breakable,colback=teal!4,colframe=teal!35!black,title={Before you move on}]
What does \kn{ನಿಲ್ಲಿಸಿ} mean? (\textquotedblleft{}Please stop\textquotedblright{}.) Say it once more.
\end{tcolorbox}
\section[\texorpdfstring{\kn{ಕರೆಯಿರಿ} (kareyiri)}{ಕರೆಯಿರಿ (kareyiri)}]{\kn{ಕರೆಯಿರಿ} (kareyiri) — please call}

\label{lesson:KA-C101-call}



\begin{quote}
Before the new one: say the Kannada for please close, then the Kannada for please stop.
\end{quote}

\subsection*{You'll want to know: \kn{ಕರೆಯಿರಿ}}
\textbf{\kn{ಕರೆಯಿರಿ}} (\emph{kareyiri}) — \textquotedblleft{}please call (someone)\textquotedblright{}.

\textbf{\kn{ವೈದ್ಯರನ್ನು} \kn{ಕರೆಯಿರಿ}} (\emph{vaidyarannu kareyiri}) — \textquotedblleft{}please call the doctor\textquotedblright{}.

\begin{cousinweb}
\textbf{\kn{ಕರೆಯಿರಿ}} (\emph{kareyiri}) is native Dravidian, from \emph{kare}, \textquotedblleft{}to call, invite\textquotedblright{}, which DEDR 1291 lists itself, with the polite ending \emph{-iri}; the entry also has Tamil \emph{karai}, Telugu \emph{kraṅgu} and Malayalam \emph{karayuka}.
\end{cousinweb}

\begin{grammarlens}[title={What you've built}]
The request when help is needed.
\end{grammarlens}

\subsection*{Your turn}
\begin{itemize}
\raggedright
  \item \emph{Say it:} \emph{kareyiri}
  \item \emph{Say it:} \emph{kareyiri}, once more
  \item \emph{Say it:} say \emph{vaidyarannu kareyiri}
  \item \emph{Recall:} say the Kannada for please close, then the Kannada for please stop, then say \emph{kareyiri} again
\end{itemize}

\begin{tcolorbox}[breakable,colback=teal!4,colframe=teal!35!black,title={Before you move on}]
What does \kn{ಕರೆಯಿರಿ} mean? (\textquotedblleft{}Please call\textquotedblright{}.) Say it once more.
\end{tcolorbox}
\section[\texorpdfstring{\kn{ತುರ್ತು} (turtu)}{ತುರ್ತು (turtu)}]{\kn{ತುರ್ತು} (turtu) — urgent, an emergency}

\label{lesson:KA-C101-urgent}



\begin{quote}
Before the new one: say the Kannada for please stop, then the Kannada for please call.
\end{quote}

\subsection*{You'll want to know: \kn{ತುರ್ತು}}
\textbf{\kn{ತುರ್ತು}} (\emph{turtu}) — \textquotedblleft{}urgent, an emergency\textquotedblright{}.

\textbf{\kn{ತುರ್ತು}!} — \textquotedblleft{}it's urgent!\textquotedblright{}

\begin{cousinweb}
\textbf{\kn{ತುರ್ತು}} (\emph{turtu}) is a loan from Persian \emph{turt}, \textquotedblleft{}quick, prompt\textquotedblright{}, through Hindustani. What must be done quickly is urgent.
\end{cousinweb}

\begin{grammarlens}[title={What you've built}]
The word that makes people hurry.
\end{grammarlens}

\subsection*{Your turn}
\begin{itemize}
\raggedright
  \item \emph{Say it:} \emph{turtu}
  \item \emph{Say it:} \emph{turtu}, once more
  \item \emph{Say it:} say \emph{turtu!}
  \item \emph{Recall:} say the Kannada for please stop, then the Kannada for please call, then say \emph{turtu} again
\end{itemize}

\begin{tcolorbox}[breakable,colback=teal!4,colframe=teal!35!black,title={Before you move on}]
What does \kn{ತುರ್ತು} mean? (\textquotedblleft{}Urgent, an emergency\textquotedblright{}.) Say it once more.
\end{tcolorbox}
